\addtocontents{toc}{\protect\Needspace{3\baselineskip}}
\originalchapter{边界条件与经典弦运动}\label{ch:classical}
\originalsection{端点与 D 膜}
对 $X$ 变分, 除波动方程外还有边界项 $\delta S\vert_{\partial\Sigma}=-(2\pi\alpha')^{-1}\int\dd\tau X'\cdot\delta X$. 自由端点要求 $X'=0$, 即 Neumann 条件. 固定方向要求 $\delta X=0$, 即 Dirichlet 条件. 对时间及 $p$ 个空间方向取 Neumann, 其余方向取 Dirichlet, 给出 D$p$ 膜的边界条件. 两端可落在不同膜上, 但每个端点的边界项都须分别消失.

在静态规范和均匀能量参数中, $\boldsymbol X(t,\sigma)=\frac12[\boldsymbol F(ct+\sigma)+\boldsymbol G(ct-\sigma)]$. 约束为 $|\boldsymbol F'|=|\boldsymbol G'|=1$. 两个函数的导数是单位球上的路径, 不是任意空间曲线. 自由开弦在 $\sigma=0$ 给 $\boldsymbol F'=\boldsymbol G'$, 在另一端给其导数周期 $2\sigma_1$. 自由端点处 $\boldsymbol X'=0$, 约束使端点速率为 $c$. 膜限制端点方向时, 端点沿膜速度可低于 $c$, 因它还受垂直张力约束.

\originalsection{旋转弦与 Regge 关系}
取单位 $c=1$, 端点自由的平面旋转解可写 $X^1=A\cos\sigma\cos\tau$, $X^2=A\cos\sigma\sin\tau$, $X^0=A\tau$, $0\le\sigma\le\pi$. 导数直接满足波动方程与式\eqref{eq:constraints}, 两端 Neumann 条件也成立. Noether 能量 $E=T_s\int_0^\pi\dot X^0\dd\sigma=\pi T_sA$, 角动量
$J=T_s\int(X^1\dot X^2-X^2\dot X^1)\dd\sigma=\pi T_sA^2/2$.
消去 $A$ 得 $J=\alpha'E^2$. 这不是任意伸展物体的普遍关系, 而是这一最简单自由旋转弦解的结果, 它说明延展物体的角动量随能量平方增长.

闭弦的左右行波各自满足周期条件. 当 $\boldsymbol F'(u)=\boldsymbol G'(v)$, 空间切向导数 $\boldsymbol X'=0$, 瞬时速度达到 $c$, 可形成 cusp. 局部形状由二阶与三阶展开控制, 与一般平滑段的近似不同. 若左右路径没有相交, 不能宣称一定存在 cusp.

\originalsection{守恒量}
平直目标空间平移给世界面动量流 $P^a_\mu=-(2\pi\alpha')^{-1}\partial^aX_\mu$. 场方程意味着 $\partial_aP^a_\mu=0$. 空间积分 $p_\mu=\int P^\tau_\mu\dd\sigma$ 的时间变化等于两端动量流差; Neumann 条件使自由方向总动量守恒, Dirichlet 固定背景可吸收相应动量.

Lorentz 变换 $\delta X^\mu=\omega^\mu{}_{\nu}X^\nu$ 给 $J^a_{\mu\nu}=X_\mu P^a_\nu-X_\nu P^a_\mu$. 它的守恒同时用到场方程与 $\partial_aX_\mu\partial^aX_\nu$ 的对称性. 显式边界条件若破坏某些目标空间对称, 相应总荷不必守恒. 这与「世界面坐标重命名」的规范对称不同, 不能把两者的 Noether 荷混在一起.

\originalsection{端点流与弦长参数}
均匀能量参数的空间约束可直接由静态规范得出. 令 $X^0=ct$, 并重新选 $\sigma$ 使每参数长度所携能量为 $T_0$. 则 $\dot{\boldsymbol X}\cdot\boldsymbol X'=0$ 以及 $|\boldsymbol X'|^2+|\dot{\boldsymbol X}|^2/c^2=1$. 前式表示所选时间方向与参数切向正交, 后式说明弦运动越快, 同样能量参数的实际空间长度越短. 总能 $E=T_0\sigma_1$ 由规范定义得到, 并不意味着所有时刻实际弧长都等于 $\sigma_1$.

自由端点 $\boldsymbol X'=0$ 时必有 $|\dot{\boldsymbol X}|=c$. 当端点在 D 膜上, 只有切向动量流为零, 法向 $X'$ 不必零, 于是端点切向速度可以更小. 例如平行膜间静直弦 $\boldsymbol X'=\hat z$, $\dot X=0$ 完全符合约束, 不与自由端点光速结论冲突. 保留边界条件的具体方向才会区分这两种情形.
